% !PRINTED 190
\weberpage{190}

Für die Factoren von $f(x, y)$, die uns jetzt öfter begegnen, führen wir als Abkürzung das Zeichen ein

\weq{6}{x\beta_1 - y\alpha_1 = (x\beta_1).}

Ebenso werden wir setzen

\weq{7}{x_1 y_2 - x_2 y_1 = (x_1 y_2),}

worin $x_1, y_1$ und $x_2, y_2$ zwei beliebige Variablenpaare sind. Wenn wir die inhomogene Darstellung anwenden, so setzen wir noch kürzer

\weq{8}{\begin{aligned}
x - \alpha_1 &= (0, 1),\ x - \alpha_2 = (0, 2) \ldots x - \alpha_n = (0, n) \\
\alpha_1 - \alpha_2 &= (1, 2) \ldots \alpha_i - \alpha_k = (i, k).
\end{aligned}}

Machen wir in der Form (2) eine lineare Substitution

\weq{9}{\begin{aligned}
x &= \alpha x' + \beta y' \\
y &= \gamma x' + \delta y'
\end{aligned}}

mit der von Null verschiedenen Determinante

\weq{10}{r = \alpha\delta - \beta\gamma,}

so geht die Form $f(x, y)$ in eine Form $n^{\text{ten}}$ Grades $F(x', y')$ über, und es ist

\weq{11}{F(x', y') = a'_0 x'^n + a'_1 x'^{n-1} y' + a'_2 x'^{n-2} y'^2 + \cdots a'_n y'^n.}

Die Coëfficienten $a'_0, a'_1 \ldots a'_n$ der transformirten Form erhält man als lineare und homogene Ausdrücke in den Coëfficienten $a_0, a_1 \ldots a_n$ und als homogene Ausdrücke $n^{\text{ten}}$ Grades in den Substitutionscoëfficienten $\alpha, \beta, \gamma, \delta$. Man erhält so, wenn man $x' = 1, y' = 0$ oder $x' = 0, y' = 1$ setzt,

\weq{12}{a'_0 = f(\alpha, \gamma), \qquad a'_n = f(\beta, \delta).}

Die übrigen Coëfficienten $a'$ kann man durch die Derivirten der Function $f(x, y)$ ausdrücken, z. B.

% [sic] Im Druck steht der Factor (n - 1); die Coëfficientenform in (2) lässt ihn rechnerisch nicht erwarten.
\[
(n - 1) a'_1 = \beta f'(\alpha) + \delta f'(\gamma),
\]

was wir nicht weiter ausführen.

Wenn wir auf die beiden Variablenpaare $x, y$ und $\alpha_k, \beta_k$ gleichzeitig die Transformation (9) anwenden, so erhalten wir die linearen Factoren von $f(x, y)$ transformirt, nämlich

\weq{13}{(x'\beta'_k) = r(x\beta_k),}

wie man aus der Multiplicationsregel der Determinanten oder auch durch directes Ausrechnen findet. Diese Factoren $(x\beta_k)$ sind also auch \gesperrt{Covarianten} von $f(x, y)$, freilich aber irrationale, da sie nicht rational von den Coëfficienten von $f(x)$ abhängen. Ebenso ergiebt sich
%CONT>>
